\section{Further questions}
\label{sec:conclusion}

The completed norm-$89$ comparison illustrates two distinct sources of
rigidity. The endpoint potential makes cost nonnegative on a large regular
language. Its equality graph then forces a count relation incompatible
with a reduced genuine occurrence, except at the fixed base source.
Finite balance adds a quantitative constraint: long paths with small
completed excess must have frequency close to a recurrent zero-cost
component.

Three questions remain particularly concrete. First, one can seek a
description of endpoint potentials that persists as the Gaussian block
varies. Any such statement must retain both primitive content and the
actual source and terminal charges. The present proof exhibits one explicit potential at $89$ and supplies no
potential uniform in the block.

Second, the graph frequencies $0,1/5,1/3,1$ need not all be attainable as
limits of genuine sources with bounded excess. The exact mechanical
recognition condition is stronger than finite balance and coprime
counts. Incorporating it into the critical-component analysis could
remove some exceptional slopes or improve the constants. The current
theorem makes neither assertion.

Third, minimization on an entire Gaussian fiber requires comparing
allocations involving other prime blocks. Strict increase under one
complete flip is a useful local statement about this arithmetic orbit,
but it does not compare every primitive representation of the same norm.
The infinite family in \cref{thm:infinite-family} provides genuine
instances where this one-block penalty becomes large; understanding
how several blocks interact remains a separate problem.

Two further directions were suggested to us, and both concern the finite
graph rather than the arithmetic. The first is the decomposition used in
the proof of \cref{lem:budget}, which cuts an accepted path at its
positive-weight edges and treats each zero-weight segment separately. That
is a first-return decomposition relative to the zero-slack subgraph, and the
theory of return words gives a sharper vocabulary for such decompositions
than the condensation budget \eqref{eq:condensation-budget} does: the
budget charges every component of a longest condensation path, whereas a
return-word induction would charge only the components a given word
actually revisits. Whether that replaces the constant $K=59$ by something
close to the six-state component alone is the concrete question.

The second is the algebraic reading of finiteness recalled in
\cref{sec:introduction}. Over a finite field, a power series is algebraic
exactly when a finite automaton generates its coefficients
\cite{ChristolKamaeMendesFranceRauzy1980}, and the reflection graph of
\cref{sec:endpoint} is a finite automaton attached to an arithmetic
operation. The analogy stops there, since that equivalence is a
positive-characteristic phenomenon and our labels are ordinary integers; we
do not know whether some series formed from the costs $\Delta$ along the
graph satisfies an algebraic relation, and we have not found a formulation
in which the question is more than suggestive.

The general theorem in \cref{thm:balanced-stability} can be used outside
the Markoff setting. It takes as input a finite graph whose initial, edge
and terminal weights are nonnegative integers and whose zero-weight
cyclic components each carry a single letter frequency, and it returns an
explicit bound on the letter count of every balanced word the graph
accepts. Separating that mechanism from the arithmetic seems to us the
most transferable part of this article.
